\enquestion{9.42}{
    A melting point test of $n = 15$ samples of a binder used in manufacturing a rocket propellant resulted in $\bar{x} = \qty{68}{\degree C}$.
    Assume that the melting point is normally distributed with $\sigma = \qty{0.83}{\degree C}$.
    \begin{enumerate}
        \item[(d)] What value of $n$ would be required if we want $\beta < 0.1$ when $\mu = 65.5$?
        Assume that $\alpha = 0.01$.
    \end{enumerate}
}

\enquestion{9.47}{
    Medical researchers have developed a new artificial heart constructed primarily of titanium and plastic.
    The heart will last and operate almost indefinitely once it is implanted in the patient's body, but the battery pack needs to be recharged about every four hours.
    A random sample of $50$ battery packs is selected and subjected to a life test.
    The average life of these batteries is $\qty{4.05}{\hour}$.
    Assume that battery life is normally distributed with standard deviation $\sigma = \qty{0.2}{\hour}$.
    \begin{enumerate}
        \item[(d)] What sample size would be required to detect a true mean battery life of $\qty{4.5}{\hour}$ if you wanted the power of the test to be at least $0.9$?
        \item[(e)] Is there evidence to support the claim that mean battery life exceeds $4$ hours? 
        Explain how this question could be answered by constructing a one-sided confidence bound on the mean life.
    \end{enumerate}
}

\enquestion{9.128}{
    Consider the following computer output.

    \begin{center}
        \begin{tabular}{*{7}{C{.11428571428\textwidth}}}
            \toprule
                \multicolumn{7}{l}{\textbf{One-Sample $T$:}} \\                
                \multicolumn{7}{l}{Test of $\mu = 85$ vs $\mu<85$} \\
                Variable & $n$ & Mean & StDev & SE Mean & $Z$ & $p$ \\
                $X$ & $31$ & $84.331$ & ? & $0.631$ & ? & ? \\
            \bottomrule
        \end{tabular}
    \end{center}

    \begin{enumerate}[label=(\alph*)]
        \item How many degrees of freedom are there on the $t$-statistic?
        \item Fill in the missing information. You may use bounds on the $P$-value.
        \item What are your conclusions if $\alpha = 0.05$?
        \item Find a \SI{95}{\percent} upper-confidence bound on the mean.
        \item What are your conclusions if the hypothesis is $H_{0}: \mu = 100$ versus $H_{0}: \mu > 100$?
    \end{enumerate}
}

\enquestion{9.130}{
    An article in \emph{Fire Technology} [\enquote{An Experimental Examination of Dead Air Space for Smoke Alarms} (2009, Vol.45,pp.97-115)] studied the performance of smoke detectors installed not less than $\qty{100}{\mm}$ from any adjoining wall if mounted on a flat ceiling, and not closer than $\qty{100}{\mm}$ and not farther than $\qty{300}{\mm}$ from the adjoining ceiling surface if mounted on walls.
    The purpose of this rule is to avoid installation of smoke alarms in the \enquote{dead air space,} where it is assumed to be difficult for smoke to reach.
    The paper described a number of interesting experiments.
    Results on the time to signal (in seconds) for one such experiment with pine stick fuel in an open bedroom using photoelectric smoke alarms are as follows: $220$, $225$, $297$, $315$, $282$, and $313$.
    \begin{enumerate}[label=(\alph*)]
        \item Is there sufficient evidence to support a claim that the mean time to signal is less than $285$ seconds?
        \item Is there practical concern about the assumption of a normal distribution as a model for the time-to-signal data?
        \item Find a \SI{95}{\percent} two-sided CI on the mean time to signal.
    \end{enumerate}
}

\enquestion{9.131}{
    Suppose we wish to test the hypothesis $H_{0}: \mu = 85$ versus the alternative $H_{1}: \mu > 85$ where $\sigma = 16$.
    Suppose that the true mean is $\mu = 86$ and that in the practical context of the problem this is not a departure from $\mu_{0} = 85$ that has practical significance.
    \begin{enumerate}[label=(\alph*)]
        \item For a test with $\alpha = 0.01$, compute $\beta$ for the sample sizes $n = 30$, $100$, $400$, and $2550$ assuming that $\mu = 86$.
        \item Suppose that the sample average is $\bar{x} = 86$.
        Find the $p$-value for the test statistic for the different sample sizes specified in part (a).
        Would the data be statistically significant at $\alpha = 0.01$?
        \item Comment on the use of a large sample size in this problem.
    \end{enumerate}
}


